\section{Diseño de la simulación}\label{sec:diseno}

La simulación usa un generador con $\theta$ conocida y cambia un factor cada vez sobre un escenario base, de modo que cada resultado teórico tiene un escenario que lo contrasta con una predicción numérica previa.

\paragraph{Generador.}
\begin{enumerate}
  \item \emph{Covariables:} $V=(V_1,\dots,V_4)$ normal con correlaciones $0,3^{|j-k|}$, más $V_5$ categórica de 3 niveles cuya distribución depende de $V_1$; todas son características del cliente. Una variable $U$ no observada, independiente de $V$, solo entra en el escenario de confusión no observada.
  \item \emph{Variable de canal latente:} $H=\ind\{C^*\ne\vacio\}$ indica historial real, con $\PP(H=1\mid V)=\sigma(-0,4+0,8V_1+0,3V_3)$; si $H=1$, el historial vale $N(0,5V_1;\,1)$. Se registra si $D=1$ y $H=1$, y $C=\vacio$ en otro caso.
  \item \emph{Canal:} $\PP(D=1\mid V,H,U)=\sigma(\alpha_0+\omega\,\alpha^\top V+\zeta H+\nu U)$. El parámetro $\omega$ controla el solapamiento, y $\zeta\ne0$ hace que la distribución de $C^*$ dado $V$ difiera entre canales, condición necesaria para que $C^*$ pueda confundir.
  \item \emph{Resultados:} dos productos comunes, $\PP(Y_k(d)=1\mid X^*,U)=\sigma(\beta_{k0}+\gamma_k d+\beta_k^\top V+\delta H+\nu'U+b)$ para $k=1,2$, con $\sigma$ la función logística y $b$ un efecto aleatorio de cliente, y un producto exclusivo del canal premium con $Y_3(0)=0$. El resultado que usan los estimadores es $Y=Y_1+Y_2$.
  \item \emph{Grupo objetivo:} $S=\ind\{D=1,\ s^*(V)\le c^*\}$, con $s^*$ el índice verdadero de propensión de compra en el canal premium y $c^*$ su percentil 10 en ese canal. La regla es fija, así que $\theta$ no cambia entre réplicas.
  \item \emph{Unidades y observaciones:} cada cliente genera entre 1 y 5 llamadas, en número independiente de todas las demás variables, con el canal y las covariables constantes entre ellas; por eso los estimandos por unidad y por observación coinciden.
  \item \emph{Verdad:} $\theta$ y los límites teóricos de cada estimador se calculan en una población de 1,5 millones de clientes; las medias condicionales verdaderas y $dP_V/dQ_V$ se obtienen integrando $H$, $U$ y el efecto de cliente por cuadratura de Gauss--Hermite. Su error de Monte Carlo, de unos 0,001 puntos (0,006 con confusión a través del historial y 0,015 con confusión no observada), es despreciable frente al de las medias de las réplicas. Los efectos del recorte, mucho más pequeños, se predicen como diferencias pareadas dentro de la población.
\end{enumerate}

\paragraph{Escenarios y predicciones.} El escenario base no incluye $C$, tiene $\zeta\ne0$ y $\delta=\nu=\nu'=0$, solapamiento alto ($\omega=0,5$; $\chi^2=3,6$), LightGBM \citep{ke2017} con \emph{cross-fitting} de 5 bloques por cliente para $\hat p$ y $\hat m$, y 30\,000 clientes por réplica: unas 62\,000 llamadas donantes y 2\,800 del grupo objetivo. Sobre él variamos la inclusión de $C$, la confusión por el historial ($\delta=1,2$), el solapamiento ($\omega=1,5$ y $3$; $\chi^2=6,3$ y unos 36; este último es inestable porque $w^2$ tiene una cola muy pesada: en seis poblaciones de 1,5 millones de clientes varió entre 33 y 63, así que el $\neff/n_0$ teórico con solapamiento bajo, 0,03, es solo orientativo), la especificación de los modelos, la confusión no observada ($\nu=\nu'=1$), el recorte en los percentiles 99, 95 y 90 de los pesos donantes y el tamaño (3\,000 a 30\,000 clientes). Cada escenario tiene 200 réplicas (las ejecuciones del bootstrap y del contraste de signo del recorte, menos; sección~\ref{sec:resultados}), con lo que la cobertura al 95\,\% se estima con un error típico de unos 1,5 puntos. Las predicciones que no se derivan directamente de la teoría se registraron con fecha antes de ejecutar el escenario correspondiente; el código, las predicciones con la fecha de su commit y los resultados brutos están disponibles en \url{https://github.com/pablotgp/channel-measurement}. Registrar significa aquí hacer un commit fechado en ese repositorio antes de la ejecución, no usar un registro externo.

\paragraph{El AUC no mide el solapamiento.} Al fijar $\omega$, el AUC del clasificador ideal se mantuvo entre 0,89 y 0,93 para $\omega$ entre 0,3 y 2, mientras $\chi^2$ variaba entre 3,5 y unos 11, sin crecer de forma monótona (es mínima para $\omega$ entre 0,5 y 1). La propia definición del grupo objetivo, un corte en una puntuación, ya lo separa del canal donante. Por eso fijamos los niveles por $\chi^2$, que es lo que gobierna la varianza (Proposición~\ref{prop:varianza}).
